\section{算术障碍与阈值下界}\label{sec:obstruction}
\begin{lemma}\label{lem:affine}
任一平均词 $W$ 满足 $W(a\mathbf1+bx)=a\mathbf1+bWx$。因此配置达到均值，等价于其本原中心化代表达到零。
\end{lemma}
\begin{proof}
每个 $T_S$ 线性且固定 $\mathbf1$，对词长归纳即可。
\end{proof}

\begin{lemma}\label{lem:obstruction}
非零本原整数零和配置满足 $G\mid n$。若素数 $\ell\ne p$ 整除 $G$，则配置不可归零。
\end{lemma}
\begin{proof}
模 $G$ 的全部坐标同余于某个 $a$。本原性给出 $\gcd(a,G)=1$，零和给出 $na\equiv0\pmod G$，所以 $G\mid n$。若 $\ell\mid G$，共同余数模 $\ell$ 非零。所有中间值属于 $\Z[1/p]$，该环可约化至 $\F_\ell$；每次平均保持这个非零共同余数，故不能归零。
\end{proof}

称一个整数零和配置是\emph{合法的}，如果它在每个素数 $\ell\mid n$、$\ell\ne p$ 处不全同余。这个定义在未取本原代表时比 $G=p^e$ 更强；从本原输入出发二者相同。称保持合法性的操作为\emph{安全操作}。记
\[
d_p(n)=\#\{\ell:\ell\mid n,\ \ell\ne p\text{ 为素数}\}.
\]

\begin{lemma}\label{lem:fraction}
设 $p$ 为奇素数，$p\le n\le2p$。整数零和配置的一条成功平均路径不可能首次产生非整数。
\end{lemma}
\begin{proof}
首次非整数平均后，最低 $p$ 进赋值 $h<0$ 恰由一个等值的 $p$ 块承担。若下一选择与此块交于 $k$ 个位置，$1\le k<p$，则其和的赋值为 $h$，新均值的赋值为 $h-1$，并形成新的最低 $p$ 块。选整个最低块时操作恒等。若选择与之不交，则只能 $n=2p$ 且选中补集；零和迫使结果成为 $(a^p,(-a)^p)$。按引理\ref{lem:affine}，它等价于 $(1^p,(-1)^p)$，被模 $2$ 的共同非零余数阻碍。故最低负赋值既不能消失，也不能通过不交操作进入可解配置。
\end{proof}

\begin{proposition}\label{prop:lower}
对奇素数 $p$，配置 $(-p,0^{p-1},1^p)$ 本原、零和、$G=1$，但不可归零。
\end{proposition}
\begin{proof}
由引理\ref{lem:fraction}，成功路径必须始终采用整数平均。若一次选择包含 $j$ 份 $1$ 和 $e$ 份 $-p$，整性要求 $j-ep\equiv0\pmod p$，故 $j=0$ 或 $p$。唯一非恒等可能是选择 $(-p,0^{p-1})$，得到 $((-1)^p,1^p)$；引理\ref{lem:obstruction}排除该配置。
\end{proof}

二平均使用\cite{mixing}引言的完全混合刻画定理。其条件 MC 是：对每个正奇数 $b$，若整数配置模 $b$ 全同余，则均值也属于这个余数类。代入本原零和 $X$ 后，MC 等价于 $G$ 无奇素因子，故给出 $n\ge4$ 的判据。混合图按拓扑顺序执行时，每个二入二出节点可用两个原位置实现。另一方面，$(-3,1,2)$ 的 $G=1$，任意首次二平均产生 $(a,a,-2a)$、$a\ne0$；后继非恒等操作只把参数变成 $-a/2$，故三元不可归零。这证明推论\ref{cor:threshold}的二平均部分。

